% -*- coding: utf-8 -*-
% !TEX program = xelatex
\documentclass[plain,noanswer,medmath]{../../template/jnuexam}
\usepackage{../../template/kysx}

\begin{document}


\examtitle{1994}{1}


\makepart{填空题}{本题共5个小题，每小题3分，满分15分}


\begin{problem}
$\lim_{x \to 0} \cot x\left(\frac{1}{\sin x} - \frac{1}{x} \right) =$ \fillin{}．
\end{problem}


\begin{problem}
曲面$z - \e^z + 2xy = 3$在点$(1,2,0)$处的切平面方程为 \fillin{}．
\end{problem}


\begin{problem}
设$u = \e^{-x}\sin \frac{x}{y}$，则$\frac{\pd^2 u}{\pdx\pd y}$
在点$\left(2,\frac{1}{\pi} \right)$处的值为 \fillin{}．
\end{problem}


\begin{problem}
设区域$D$为$x^2 + y^2 \le R^2$，则$\iint_D \left(\frac{x^2}{a^2} + \frac{y^2}{b^2} \right)\dx\dy =$ \fillin{}．
\end{problem}


\begin{problem}
已知$\alpha = (1,2,3),\beta = \left(1,\frac{1}{2},\frac{1}{3}\right)$，设$A = \alpha^T\beta$,
其中$\alpha^T$是$\alpha$的转置，则$A^n =$ \fillin{}．
\end{problem}


\makepart{选择题}{本题共5个小题，每小题3分，满分15分}


\begin{problem}
设$M = \int_{-\frac{\pi}{2}}^{\frac{\pi}{2}} \!\frac{\sin x}{1 + x^2}\cos^4 x\dx$, 
  $N = \int_{-\frac{\pi}{2}}^{\frac{\pi}{2}} \!(\sin^3 x + \cos^4 x)\dx$, \goodbreak 
  $P = \int_{-\frac{\pi}{2}}^{\frac{\pi}{2}} \!(x^2\sin^3 x - \cos^4 x)\dx$，则有\pickout{}
\begin{abcd}
\item $N < P < M$．                        
\item $M < P < N$．
\item $N < M < P$．                        
\item $P < M < N$．
\end{abcd}
\end{problem}


\begin{problem}
二元函数$f(x,y)$在点$(x_0,y_0)$处两个偏导数$f_x'(x_0,y_0)$,$f_y'(x_0,y_0)$存在是$f(x,y)$在该点连续的\pickout{}
\begin{abcd}
\item 充分条件但非必要条件．         
\item 必要条件而非充分条件．
\item 充分必要条件．                 
\item 既非充分条件又非必要条件．
\end{abcd}
\end{problem}


\begin{problem}
设常数$\lambda > 0$，且级数$\sum_{n = 1}^{\infty}a_n^2$收敛，则级数$\sum_{n = 1}^{\infty}(- 1)^n\frac{|a_n|}{\sqrt{n^2 + \lambda}}$\pickout{}
\begin{abcd}
\item 发散．                               
\item 条件收敛．
\item 绝对收敛．                           
\item 收敛性与$\lambda$有关．
\end{abcd}
\end{problem}


\begin{problem}
设$\lim_{x \to 0} \frac{a\tan x + b(1 - \cos x)}{c\ln(1 - 2x) + d(1 - \e^{- x^2})} = 2$，
其中$a^2 + c^2 \ne 0$，则必有\pickout{}
\begin{abcd}
\item $b = 4d$．    
\item $b = - 4d$．    
\item $a = 4c$．    
\item $a = - 4c$．
\end{abcd}
\end{problem}


\begin{problem}
已知向量组$\alpha_1$, $\alpha_2$, $\alpha_3$, $\alpha_4$线性无关，则向量组\pickout{}
\begin{abcd}
\item $\alpha_1 + \alpha_2$、$\alpha_2 + \alpha_3$、$\alpha_3 + \alpha_4$、$\alpha_4 + \alpha_1$线性无关．
\item $\alpha_1 - \alpha_2$、$\alpha_2 - \alpha_3$、$\alpha_3 - \alpha_4$、$\alpha_4 - \alpha_1$线性无关．
\item $\alpha_1 + \alpha_2$、$\alpha_2 + \alpha_3$、$\alpha_3 + \alpha_4$、$\alpha_4 - \alpha_1$线性无关．
\item $\alpha_1 + \alpha_2$、$\alpha_2 + \alpha_3$、$\alpha_3 - \alpha_4$、$\alpha_4 - \alpha_1$线性无关．
\end{abcd}
\end{problem}


\makepart{计算题}{本题共3小题，每小题5分，满分15分}


\begin{problem}
设$\begin{cases}
 x = \cos(t^2), \\
 y = t\cos(t^2) - \int_1^{t^2} \frac{1}{2\sqrt{u}}\cos u\du,
\end{cases}$ 求$\frac{\dy}{\dx}$、$\frac{\d^2y}{\dx^2}$在$t = \sqrt{\frac{\pi}{2}}$的值．
\end{problem}


\begin{problem}
将函数$f(x) = \frac{1}{4}\ln \frac{1 + x}{1 - x} + \frac{1}{2}\arctan x - x$展开成$x$的幂级数．
\end{problem}


\begin{problem}
求$\int \frac{\dx}{\sin 2x + 2\sin x}$．
\end{problem}


\makepart{}{本题满分6分}

\begin{problem}
计算曲面积分$\iint_S \frac{x\dy\dz + z^2 \dx\dy}{x^2 + y^2 + z^2}$,
其中$S$是由曲面$x^2 + y^2 = R^2$及两平面$z = R$, $z = - R$($R > 0$)所围成立体表面的外侧．
\end{problem}


\makepart{}{本题满分9分}

\begin{problem}
设$f(x)$具有二阶连续导数，$f(0) = 0,f'(0) = 1$，且
$$[xy(x + y) - f(x)y]\dx + [f'(x) + x^2 y]\dy = 0$$
为一全微分方程，求$f(x)$及此全微分方程的通解．
\end{problem}


\makepart{}{本题满分8分}

\begin{problem}
设$f(x)$在点$x = 0$的某一邻域内具有二阶连续导数，且$\lim_{x \to 0} \frac{f(x)}{x} = 0$，证明级数
$\sum_{n = 1}^{\infty}f\left(\frac{1}{n} \right)$绝对收敛．
\end{problem}


\makepart{}{本题满分6分}

\begin{problem}
已知点$A$与$B$的直角坐标分别为$(1,0,0)$与$(0,1,1)$．线段$AB$绕$z$轴旋转一周所围成的旋转曲面为$S$．求由$S$及两平面$z = 0,z = 1$所围成的立体体积．
\end{problem}


\makepart{}{本题满分8分}

\begin{problem}
设四元线性齐次方程组①为$\begin{cases}
  x_1 + x_2 = 0, \\
  x_2 - x_4 = 0,
\end{cases}$ 又已知某线性齐次方程组②的通解为
$k_1(0,1,10) + k_2(- 1,2,2,1)$．
\begin{enumerate}
\item 求线性方程组①的基础解系；
\item 问线性方程组①和②是否有非零公共解？
    若有，则求出所有的非零公共解．若没有，则说明理由．
\end{enumerate}
\end{problem}


\makepart{}{本题满分6分}

\begin{problem}
设$A$为$n$阶非零方阵，$A^*$是$A$的伴随矩阵，$A^T$是$A$的转置矩阵，
当$A^* = A^T$时，证明$|A| \ne 0$．
\end{problem}


\makepart{填空题}{本题共2小题，每小题3分，满分6分}


\begin{problem}
已知$A$,$B$两个事件满足条件$P(AB) = P(\overline{A}\overline{B})$，且$P(A) = p$，则$P(B) =$ \fillin{}．
\end{problem}


\begin{problem}
设相互独立的两个随机变量$X$、$Y$具有同一分布律，且$X$的分布律为
$$\begin{array}{|c|c5c|}
\hline
  X &           0 & 1 \\
\hline
  P & \frac{1}{2} & \frac{1}{2} \\
\hline
\end{array}$$
则随机变量$Z = \max \left\{X,Y \right\}$的分布律为 \fillin{}．
\end{problem}


\makepart{}{本题满分6分}

\begin{problem}
已知随机变量$(X,Y)$服从二维正态分布，且$X$和$Y$分别服从正态分布$N(1,3^2)$和
$N(0,4^2)$, $X$与$Y$的相关系数$\rho_{XY} = - \frac{1}{2}$，设$Z = \frac{X}{3} + \frac{Y}{2}$．
\begin{enumerate}
\item 求$Z$的数学期望$E(Z)$和方差$D(Z)$；
\item 求$X$与$Z$的相关系数$\rho_{XZ}$；
\item 问$X$与$Z$是否相互独立？为什么？
\end{enumerate}
\end{problem}


\end{document}
